\section{Exact finite screening and fixed-background transitions}
\label{sec:screening}

We next determine an exact small-displacement coefficient for a symmetric
finite screen and show that it survives every fixed polynomial-count axis
background. The background may define an unbounded positive form on a
fixed window. The argument does not assume its boundedness and does not
let either the screen size or the background vary with the displacement.
The finite-dimensional flat-limit and Schur-complement tools used below
are established perturbation methods; see \cite{UsevichBarthelme}.
We give the signed calculation and the smooth-background transfer in full.

\subsection{The coefficient and the two asymptotic conclusions}

For an integer \(n\ge1\), define the divisor
\begin{equation}\label{scr:screen-definition}
 \mathcal S_{n,\beta}
 =128\sum_{j=1}^n
   \left[i\beta n\cos\frac{(2j-1)\pi}{2n}\right]
       +[\beta]+[-\beta],
 \qquad \beta>0.
\end{equation}
Here \([z]\) is an atom of unit multiplicity. Thus the screen has total
multiplicity \(128n+2\), although only \(n+2\) distinct locations.
Put
\begin{equation}\label{scr:parameters}
 r=\lfloor n/2\rfloor,\qquad d=2r+1,\qquad
 t_j=n\cos\frac{(2j-1)\pi}{2n}>0\quad(1\le j\le r),
 \qquad \alpha_n=\operatorname{arsinh}(1/n).
\end{equation}
An empty sum is zero and an empty product is one. Define
\begin{align}
 \kappa_n
 &=\frac4n\sum_{k=0}^{r-1}
              \sinh^2((2k+1)\alpha_n)
   =\frac1n\left\{
       \frac{\sinh(4r\alpha_n)}{\sinh(2\alpha_n)}-2r
            \right\},                                      \label{scr:kappa}\\
 P_n
 &=\prod_{j=1}^r(1+t_j^2)
   =n^n2^{1-n}
       \begin{cases}
        \cosh(n\alpha_n),&n\ \text{even},\\
        \sinh(n\alpha_n),&n\ \text{odd},
       \end{cases}                                         \label{scr:product}\\
 h_j
 &=\frac{2^{2j+1}(j!)^4}{(2j+1)((2j)!)^2}.                 \label{scr:legendre-norm}
\end{align}
In particular, \(P_1=1\), \(\kappa_1=0\), and
\(0\le\kappa_n<16<128\). The positive coefficient governing the
screen is
\begin{equation}\label{scr:coefficient}
 c_n=\frac{2P_n^2h_d}{(d!)^2(1-\kappa_n/128)}
     =\frac{2^{2d+2}(d!)^2P_n^2}
       {(2d+1)((2d)!)^2(1-\kappa_n/128)}.
\end{equation}

On \(L^2(-L,L)\), the finite-screen operator has exactly one
negative eigenvalue, denoted by \(\mu_-(n,\beta,L)\). For each fixed
\(n\), with \(\epsilon=\beta L\),
\begin{equation}\label{scr:eigenvalue}
 \mu_-(n,\beta,L)
   =-L c_n\epsilon^{2d}
                   \bigl(1+O_n(\epsilon^2)\bigr),
 \qquad \epsilon\longrightarrow0^+.
\end{equation}
The estimate depends on \(\beta,L\) only through the displayed factor
\(L\) and the parameter \(\epsilon\). Its error is not claimed to
be uniform in \(n\). The leading power in the magnitude of the
negative eigenvalue is
\[
 2d=
 \begin{cases}
  2n,&n\ \text{odd},\\
  2n+2,&n\ \text{even}.
 \end{cases}
\]
After unitary scaling to \((-1,1)\), a normalized negative
eigenfunction tends, up to a unit scalar, to the normalized odd
Legendre polynomial of degree \(d\).

Now let \(\mathcal A\) be a fixed locally finite divisor on
\(i\R\), with positive integer multiplicities and
\begin{equation}\label{scr:background-count}
 N_{\mathcal A}(Y):=
   \sum_{\substack{i\gamma\in\mathcal A\\|\gamma|\le Y}}
      \nu_{\mathcal A}(i\gamma)
 \le C_{\mathcal A}(1+Y)^{m_{\mathcal A}},
 \qquad Y\ge0,
\end{equation}
for finite constants \(C_{\mathcal A}\ge0\) and
\(m_{\mathcal A}\ge0\). The divisor may be empty, may contain the
origin, and need not be separated at large ordinates. Set
\[
 Z_\beta=\mathcal A\sqcup\mathcal S_{n,\beta},\qquad
 q_n=\frac{2d}{2d+1},
\]
where the union adds multiplicities, and define
\begin{equation}\label{scr:margin-definition}
 \cM_{Z_\beta}(L)
 =\sup_{\substack{0\ne f\in C_c^\infty((-L,L))}}
       \max\left\{0,-\frac{Q_{Z_\beta}(f,f)}{\norm f_2^2}\right\}.
\end{equation}
For fixed \(n\) and fixed \(\mathcal A\),
\begin{equation}\label{scr:transition}
 \cM_{Z_\beta}(x\beta^{-q_n})
       \longrightarrow c_nx^{2d+1}
 \quad\text{locally uniformly for }x\in(0,\infty),
 \qquad \beta\longrightarrow0^+.
\end{equation}
In particular, for each fixed \(\eta>0\), let
\[
 L_\eta(\beta)
 =\inf\left\{L>0:\ \exists\,0\ne f\in C_c^\infty((-L,L)),
           \ Q_{Z_\beta}(f,f)\le-\eta\norm f_2^2\right\},
\]
with the infimum of the empty set equal to \(+\infty\). Then
\begin{equation}\label{scr:first-radius}
 L_\eta(\beta)\sim
       \left(\frac{\eta}{c_n}\right)^{1/(2d+1)}
                     \beta^{-2d/(2d+1)}.
\end{equation}
The same asymptotic holds if the defining inequality is strict.
Thus neither separation nor the size of the fixed axis background
enters the leading constant. No uniformity over varying backgrounds
is asserted.

\subsection{Scaling, parity, and exact inertia}

The unitary change of variables
\[
 g(x)=\sqrt L\,f(Lx),\qquad -1<x<1,
\]
preserves the norm and satisfies
\(H_f(z)=\sqrt L\,H_g(Lz)\). Consequently the screen form on
\((-L,L)\) is \(L\) times the screen form on \((-1,1)\) with
parameter \(\epsilon=\beta L\). We work on \((-1,1)\) until
scaling back at the end.

The nonzero imaginary nodes occur in the pairs
\(\pm i\epsilon t_j\), with an additional zero node if \(n\)
is odd. Pairing the two imaginary nodes expresses their contribution
as a positive cosine square plus a positive sine square. The real
pair contributes
\[
 2\left|\langle g,\cosh(\epsilon x)\rangle\right|^2
       -2\left|\langle g,\sinh(\epsilon x)\rangle\right|^2.
\]
All features here are real. The even and odd subspaces are therefore
orthogonal invariant subspaces. The even block is a sum of positive
squares, including the constant square if there is a zero node. On
the odd subspace the form is exactly
\begin{equation}\label{scr:odd-block}
 q_\epsilon(g)
   =256\sum_{j=1}^r
       |\langle g,\sin(\epsilon t_jx)\rangle|^2
       -2|\langle g,\sinh(\epsilon x)\rangle|^2.
\end{equation}

For every \(\epsilon>0\), the \(r+1\) odd features in
\eqref{scr:odd-block} are linearly independent. Indeed their entire
continuations are combinations of exponentials with distinct
exponents \(\pm i\epsilon t_j,\pm\epsilon\), and distinct
exponentials are independent: differentiating an identically zero
combination at zero through one less than the number of exponents
gives an invertible Vandermonde system. The corresponding even
features are independent by the same argument, adding exponent zero
when necessary.

For completeness, if \(\Phi\) is the synthesis map for finitely
many independent features and \(J\) is their diagonal matrix of
nonzero signed weights, then its operator is \(\Phi J\Phi^*\).
Writing \(\Phi=ER\), where \(E\) is an isometry onto its range
and \(R\) is invertible, shows that its nonzero inertia is that
of \(RJR^*\), hence that of \(J\). Applied to the odd block
this gives \(r\) positive eigenvalues and one negative eigenvalue.
The even block has \(r+1\) positive eigenvalues for even \(n\),
and \(r+2\) for odd \(n\). Thus the full screen has rank \(n+2\),
positive inertia \(n+1\), negative inertia one, and an
infinite-dimensional nullspace.

\subsection{The signed moment matrix and its Schur complement}

Write \(d_k=2k+1\) for \(0\le k\le r\). The feature series are
\[
 \sin(\epsilon t_jx)
  =\sum_{k\ge0}(-1)^kt_j^{2k+1}
             \frac{\epsilon^{2k+1}x^{2k+1}}{(2k+1)!},
 \qquad
 \sinh(\epsilon x)
  =\sum_{k\ge0}
             \frac{\epsilon^{2k+1}x^{2k+1}}{(2k+1)!}.
\]
Define the real \((r+1)\)-square matrix \(F\), with columns
indexed by \(0,\ldots,r\), by
\[
 F_{j,k}=(-1)^kt_j^{2k+1}\quad(1\le j\le r),\qquad
 F_{r+1,k}=1.
\]
Let \(J_s=\operatorname{diag}(256,\ldots,256,-2)\) and
\(A=F^TJ_sF\). Its entries are
\begin{equation}\label{scr:moment-matrix}
 A_{k\ell}
  =256(-1)^{k+\ell}\sum_{j=1}^r t_j^{2(k+\ell+1)}-2.
\end{equation}
Factoring \(t_j\) from the first \(r\) rows of \(F\) leaves
the Vandermonde matrix at \(-t_1^2,\ldots,-t_r^2,1\).
Thus \(F\) is invertible and \(A\) has inertia \((r,1)\).

For \(r\ge1\), let
\[
 V_{j,k}=(-1)^kt_j^{2k+1}
       \quad(1\le j\le r,\ 0\le k<r),\qquad
 B=V^TV,\qquad {\bf1}=(1,\ldots,1)^T.
\]
The matrix \(V\) is invertible by the same Vandermonde argument,
so \(B\) is positive definite. The leading \(r\)-square block
of \(A\) is
\[
 A_{<r}=256B-2{\bf1}{\bf1}^T,
\]
and it is positive definite precisely when
\({\bf1}^TB^{-1}{\bf1}<128\). We next verify that
\[
 {\bf1}^TB^{-1}{\bf1}=\kappa_n.
\]

For an odd polynomial of degree at most \(2r-1\), write
\[
 R(t)=\sum_{k=0}^{r-1}a_kT_{2k+1}(t/n).
\]
Discrete Chebyshev orthogonality on the original \(n\) nodes
gives
\[
 \sum_{j=1}^r|R(t_j)|^2=\frac n4\sum_{k=0}^{r-1}|a_k|^2.
\]
To see the normalization, every nonconstant Chebyshev polynomial
of degree less than \(n\) has discrete squared norm \(n/2\)
on all \(n\) nodes, with distinct degrees orthogonal, and pairing
the positive and negative nodes halves the squared norm of an odd
polynomial. A possible zero node contributes nothing.
Explicitly, with \(\theta_j=(2j-1)\pi/(2n)\), the finite geometric
sum gives \(\sum_{j=1}^n\cos(q\theta_j)=0\) for
\(1\le q\le2n-1\). Applying
\(2\cos u\cos v=\cos(u-v)+\cos(u+v)\) proves the claimed
orthogonality and normalization, since \(T_k(\cos\theta)=\cos(k\theta)\).
Also,
\[
 T_{2k+1}(i/n)=i(-1)^k\sinh((2k+1)\alpha_n).
\]
The squared norm of evaluation \(R\mapsto R(i)\), for the
positive-node norm, is therefore the first expression in
\eqref{scr:kappa}. In the alternative basis
\(R(t)=\sum_{k<r}c_k(-1)^kt^{2k+1}\), its squared node norm
is \(c^*Bc\) and \(R(i)=i{\bf1}^Tc\). The squared norm of the
same functional is \({\bf1}^TB^{-1}{\bf1}\), as required.
Summing \(2\sinh^2 v=\cosh(2v)-1\) in a finite geometric
progression gives the second expression in \eqref{scr:kappa}.
Because \(\alpha_n\le1/n\), \(2r\le n\), and
\((2r-1)\alpha_n<1\) when \(r\ge1\),
\[
 0\le\kappa_n<\frac{4r}{n}\sinh^2(1)
       \le2\sinh^2(1)<2e^2<16<128.
\]
For \(r=0\), the same conclusion holds with \(\kappa_1=0\).
This proves positive definiteness of every leading moment block
before the last one.

To compute the last Schur complement, put
\[
 a(s)=\prod_{j=1}^r(s+t_j^2)=\sum_{k=0}^r a_ks^k,
 \qquad a_r=1.
\]
Then
\[
 \sum_{k=0}^r a_k(-1)^kt_j^{2k+1}=0,\qquad
 \sum_{k=0}^r a_k=P_n.
\]
If \(c=-a_{<r}\), the last column of the positive-feature
matrix is \(Vc\). In the quadratic form of \(A\), replace
\((x,z)\in\R^r\times\R\) by \(y=x+cz\). Its value becomes
\[
 256y^TBy-2({\bf1}^Ty+P_nz)^2.
\]
The minimum over \(y\), with \(z\) fixed, has coefficient
\begin{align*}
 \sigma_n
 &=-2P_n^2
     -4P_n^2{\bf1}^T(256B-2{\bf1}{\bf1}^T)^{-1}{\bf1}\\
 &=-\frac{2P_n^2}{1-\kappa_n/128}<0.
\end{align*}
For the second equality, solving the rank-one perturbation gives
\[
 (256B-2{\bf1}{\bf1}^T)^{-1}{\bf1}
     =\frac{B^{-1}{\bf1}}{256(1-\kappa_n/128)}.
\]
Thus
\begin{equation}\label{scr:schur}
 \operatorname{Schur}(A\,|\,A_{<r})
       =\sigma_n=-\frac{2P_n^2}{1-\kappa_n/128}.
\end{equation}
When \(r=0\), this means simply \(A=(-2)\).

The product formula \eqref{scr:product} is also explicit.
For even \(n=2r\), factor
\[
 T_n(z)=2^{n-1}\prod_{j=1}^r
                    \bigl(z^2-(t_j/n)^2\bigr)
\]
and evaluate at \(z=i/n\), using
\(T_n(i/n)=(-1)^r\cosh(n\alpha_n)\).
For odd \(n=2r+1\), the corresponding factorization has the
additional factor \(z\), and
\(T_n(i/n)=i(-1)^r\sinh(n\alpha_n)\). These two evaluations
give \eqref{scr:product}, including \(n=1\).

\subsection{A nonsingular feature basis and the eigenvalue expansion}

Let \(\phi(\epsilon)\) be the column vector of the \(r\) sine
features followed by \(\sinh(\epsilon x)\), and put
\[
 D_\epsilon=\operatorname{diag}_{0\le k\le r}
                   \frac{\epsilon^{d_k}}{d_k!},\qquad
 \psi(\epsilon)=D_\epsilon^{-1}F^{-1}\phi(\epsilon).
\]
The apparent singularity at zero is removable. In fact the first
\(r+1\) odd Taylor coefficients of \(F^{-1}\phi\) vanish
except at their selected component. The remaining series converge
in \(L^2(-1,1)\), and show
\begin{equation}\label{scr:confluent-features}
 \psi_k(\epsilon,x)
      =x^{d_k}+O_{n,L^2}(\epsilon^{2(r+1-k)}).
\end{equation}
Each component extends analytically in the real parameter
\(\epsilon^2\) through zero. If \(\Psi_\epsilon\) denotes
synthesis from these columns, the identity
\(\phi=FD_\epsilon\psi\) gives the exact odd-block operator
\begin{equation}\label{scr:operator-factor}
 \Psi_\epsilon D_\epsilon A D_\epsilon\Psi_\epsilon^*.
\end{equation}

Apply Gram--Schmidt to the columns of \(\Psi_\epsilon\) in
increasing degree:
\[
 \Psi_\epsilon=E_\epsilon R_\epsilon.
\]
Here \(E_\epsilon\) has orthonormal columns, and
\(R_\epsilon\) is upper triangular with positive diagonal.
The limiting columns are independent odd monomials. The finitely
many limiting squared Gram--Schmidt norms are therefore positive;
the subtraction, division, and positive square-root operations
show that \(E_\epsilon\) and \(R_\epsilon\) are analytic in
\(\epsilon^2\) near zero.

Let \(p_j\) be the monic Legendre polynomial of degree \(j\)
on \((-1,1)\). The usual degree-\(j\) Legendre polynomial has
leading coefficient \((2j)!/(2^j(j!)^2)\) and squared norm
\(2/(2j+1)\). Dividing by the square of this leading coefficient
gives \(\norm{p_j}_2^2=h_j\) in
\eqref{scr:legendre-norm}. Gram--Schmidt on the odd monomials
consequently gives
\[
 (R_0)_{kk}=\sqrt{h_{d_k}},\qquad
 (E_0)_k=p_{d_k}/\sqrt{h_{d_k}}.
\]

Set \(V_\epsilon=D_\epsilon^{-1}R_\epsilon D_\epsilon\).
For \(j<k\), its \((j,k)\) entry is that of \(R_\epsilon\)
times \(\epsilon^{2(k-j)}d_j!/d_k!\); entries below the
diagonal vanish. Therefore
\[
 V_\epsilon=
   \operatorname{diag}(\sqrt{h_{d_0}},\ldots,\sqrt{h_{d_r}})
       +O_n(\epsilon^2).
\]
In the orthonormal coordinates \(E_\epsilon\), the operator
\eqref{scr:operator-factor} becomes
\begin{equation}\label{scr:scaled-matrix}
 M_\epsilon=D_\epsilon C_\epsilon D_\epsilon,\qquad
 C_\epsilon=V_\epsilon A V_\epsilon^T
   =\operatorname{diag}(\sqrt{h_{d_k}})\,
       A\,\operatorname{diag}(\sqrt{h_{d_k}})
           +O_n(\epsilon^2).
\end{equation}
All errors here concern fixed finite matrices at fixed \(n\).

For \(r\ge1\), write
\[
 M_\epsilon=
   \begin{pmatrix}B_\epsilon&u_\epsilon\\
                   u_\epsilon^T&a_\epsilon\end{pmatrix},
 \qquad
 S_\epsilon=a_\epsilon-u_\epsilon^TB_\epsilon^{-1}u_\epsilon.
\]
The leading block of \(C_\epsilon\), and hence \(B_\epsilon\),
is positive definite for small positive \(\epsilon\).
Taking the Schur complement in \eqref{scr:scaled-matrix} and
using \eqref{scr:schur} gives
\begin{equation}\label{scr:scaled-schur}
 S_\epsilon=
   \frac{\epsilon^{2d}}{(d!)^2}
          \{h_d\sigma_n+O_n(\epsilon^2)\}<0.
\end{equation}
Indeed diagonal scaling factors out the final diagonal square,
and the Schur complement of the limiting \(C_\epsilon\)
is \(h_d\sigma_n\).

Also
\begin{equation}\label{scr:small-elimination}
 B_\epsilon^{-1}u_\epsilon=O_n(\epsilon^2).
\end{equation}
To verify this, write \(D_<\) for the leading diagonal block of
\(D_\epsilon\) and \(c_<\) for the upper part of the last
column of \(C_\epsilon\). The vector in question is
\[
 D_<^{-1}(C_{\epsilon,<})^{-1}c_<
                  \frac{\epsilon^d}{d!}.
\]
The middle factors are uniformly bounded. Its coordinate with
index \(k<r\) is thus \(O_n(\epsilon^{d-d_k})\), and
\(d-d_k\ge2\).

Completing the square yields
\[
 (x,z)^TM_\epsilon(x,z)=
 (x+B_\epsilon^{-1}u_\epsilon z)^TB_\epsilon
       (x+B_\epsilon^{-1}u_\epsilon z)+S_\epsilon z^2.
\]
Since \(S_\epsilon<0\), this is bounded below by
\(S_\epsilon(\norm x^2+|z|^2)\). Taking
\(x=-B_\epsilon^{-1}u_\epsilon z\) gives the other Rayleigh
bound. Thus the least eigenvalue \(\lambda_-(\epsilon)\) obeys
\[
 S_\epsilon\le\lambda_-(\epsilon)\le
    \frac{S_\epsilon}
           {1+\norm{B_\epsilon^{-1}u_\epsilon}^2}
       =S_\epsilon\bigl(1+O_n(\epsilon^4)\bigr).
\]
For \(r=0\), the matrix is scalar and
\eqref{scr:scaled-schur} gives its eigenvalue directly.
Combining these estimates with \eqref{scr:schur} proves
\[
 \lambda_-(\epsilon)
  =-c_n\epsilon^{2d}\bigl(1+O_n(\epsilon^2)\bigr).
\]
Unitary scaling proves \eqref{scr:eigenvalue}.

The eigenfunction limit also follows from these bounds. The matrix is
real symmetric and its negative eigenvalue is simple, so a real
normalized negative eigenvector may be chosen. For \(r\ge1\),
such an eigenvector \((x,z)\)
satisfies
\[
 x=-(B_\epsilon-\lambda_-(\epsilon)I)^{-1}u_\epsilon z.
\]
Because \(B_\epsilon\) is positive definite and
\(\lambda_-(\epsilon)<0\), the operator
\((I-\lambda_-(\epsilon)B_\epsilon^{-1})^{-1}\)
has norm at most one. Equation \eqref{scr:small-elimination}
therefore gives \(x=O_n(\epsilon^2)|z|\).
After choosing the sign, \(z\to1\), and the limiting eigenfunction
is the last column of \(E_0\), namely \(p_d/\sqrt{h_d}\).
The same conclusion is immediate when \(r=0\).
All nonzero eigenvalues are captured by these finite coordinates.
Density of \(C_c^\infty((-L,L))\) in \(L^2(-L,L)\), and
boundedness of the finite-screen operator, show that its Rayleigh
infimum on the compact smooth class is the same negative eigenvalue.
In particular, with \(n\) fixed, the exact negative envelope is
\begin{equation}\label{scr:envelope}
 e_{n,\beta}(L):=\cM_{\mathcal S_{n,\beta}}(L)
 =-\mu_-(n,\beta,L)
 =c_n\beta^{2d}L^{2d+1}\bigl(1+O_n((\beta L)^2)\bigr)
 \quad(\beta L\longrightarrow0).
\end{equation}
For every positive \(\beta,L\) it is finite and strictly positive.
It is nondecreasing in \(L\), since the compact smooth test classes
increase with the window. The supremum need not be attained by a
compact smooth test, but such unit tests approach it arbitrarily closely.
For even \(n\), the degree is \(d=n+1\), so the powers in
\eqref{scr:envelope} are \(2n+2\) and \(2n+3\), respectively.

For reference, the first coefficients are
\begin{equation}\label{scr:first-coefficients}
 c_1=\frac43,\qquad
 c_2=\frac{1024}{44625},\qquad
 c_3=\frac{23064}{151025}.
\end{equation}
For \(n=1\) the odd block is rank one and its negative eigenvalue
is exactly
\[
 2L-\frac{\sinh(2\beta L)}{\beta}
      =-\frac43\beta^2L^3+O(\beta^4L^5).
\]
For \(n=2\), \(t_1^2=2\), \(P_2=3\), and \(\kappa_2=1/2\).
For \(n=3\), \(t_1^2=27/4\), \(P_3=31/4\), and
\(\kappa_3=4/27\). Substitution in \eqref{scr:coefficient}
gives the remaining values in \eqref{scr:first-coefficients}.

\subsection{Smooth finite-dimensional recovery of the coefficient}

To add a possibly unbounded positive axis form, we use tests from
a fixed finite-dimensional smooth space rather than approximating
an eigenfunction in the graph norm of that axis form.
Let
\[
 \mathcal V=\operatorname{span}\{x,x^3,\ldots,x^d\}
                     \subset L^2_{\rm odd}(-1,1).
\]
Choose even cutoffs \(\chi_\delta\in C_c^\infty((-1,1))\),
with \(0\le\chi_\delta\le1\), tending pointwise to one as
\(\delta\to0^+\), and put
\(\mathcal H_\delta=\chi_\delta\mathcal V\).
For sufficiently small \(\delta\), these spaces have dimension
\(r+1\). Indeed their basis vectors converge in \(L^2\) to
the independent monomials. If \(P_\delta\) denotes orthogonal
projection onto \(\mathcal H_\delta\), best approximation gives
\begin{equation}\label{scr:projected-moments}
 \norm{P_\delta x^{d_k}-x^{d_k}}_2
 \le\norm{\chi_\delta x^{d_k}-x^{d_k}}_2\longrightarrow0,
 \qquad 0\le k\le r.
\end{equation}
The projected moment vectors are consequently independent for
small \(\delta\).

Fix one such space \(\mathcal H=\mathcal H_\delta\), and let
\(P_{\mathcal H}\) be its projection. Compress
\eqref{scr:odd-block} to \(\mathcal H\). Its feature vectors
are \(P_{\mathcal H}\phi_j(\epsilon)\), and the exact
factorization \eqref{scr:operator-factor} holds with
\(\Psi_\epsilon\) replaced by \(P_{\mathcal H}\Psi_\epsilon\).
Their nonsingular limiting columns are
\(P_{\mathcal H}x^{d_k}\), by \eqref{scr:confluent-features}.
Gram--Schmidt on these columns gives a positive diagonal
\(\sqrt{h_{k,\mathcal H}}\), where \(h_{k,\mathcal H}\)
is the squared distance of \(P_{\mathcal H}x^{d_k}\) from the
span of the earlier projected moments. The triangular scaling
used in \eqref{scr:scaled-matrix} therefore gives, explicitly,
\[
 C_{\epsilon,\mathcal H}
    =\operatorname{diag}(\sqrt{h_{k,\mathcal H}})\,
          A\,\operatorname{diag}(\sqrt{h_{k,\mathcal H}})
                  +O_{n,\mathcal H}(\epsilon^2).
\]
Its leading \(r\)-square block is positive definite, and its
last Schur complement is
\(h_{r,\mathcal H}\sigma_n+O_{n,\mathcal H}(\epsilon^2)\).
The same diagonal powers yield
\(B_{\epsilon,\mathcal H}^{-1}u_{\epsilon,\mathcal H}
=O_{n,\mathcal H}(\epsilon^2)\); the completing-square
Rayleigh bounds above apply without change. The compressed
operator thus has exactly one negative eigenvalue for small
\(\epsilon>0\), equal to
\begin{equation}\label{scr:compressed-eigenvalue}
 -c_n(\mathcal H)\epsilon^{2d}
               \bigl(1+O_{n,\mathcal H}(\epsilon^2)\bigr),
 \qquad
 c_n(\mathcal H)=
   \frac{2P_n^2\rho_{\mathcal H}}
                  {(d!)^2(1-\kappa_n/128)},
\end{equation}
where
\begin{equation}\label{scr:projected-distance}
 \rho_{\mathcal H}
   =\operatorname{dist}_{L^2}^2
        \left(P_{\mathcal H}x^d,\
          \operatorname{span}\{P_{\mathcal H}x,\ldots,
                               P_{\mathcal H}x^{d-2}\}\right)>0.
\end{equation}
For \(r=0\), the span in this expression is \(\{0\}\).

Equation \eqref{scr:projected-moments} implies entrywise
convergence of the finite moment Gram matrices to their positive
definite unprojected Gram matrix. In the distance formula, the
inverse of the leading Gram block is therefore continuous, with
the scalar interpretation when \(r=0\). It follows that
\begin{equation}\label{scr:coefficient-recovery}
 \rho_{\mathcal H_\delta}\longrightarrow h_d,\qquad
 c_n(\mathcal H_\delta)\longrightarrow c_n.
\end{equation}
Moreover \(\rho_{\mathcal H}\le h_d\): project the unprojected
minimizing residual \(p_d\), use it as an admissible residual in
\eqref{scr:projected-distance}, and use the contractivity of
orthogonal projection. Thus \(0<c_n(\mathcal H)\le c_n\).

Every vector in the fixed space \(\mathcal H\) has support in
the same compact subinterval of \((-1,1)\). All norms on this
finite-dimensional space are equivalent. Consequently, for every
integer \(N\ge1\), integration by parts \(N\) times, without
boundary terms, gives a finite constant \(C_{\mathcal H,N}\)
such that
\begin{equation}\label{scr:smooth-decay}
 |H_g(it)|\le C_{\mathcal H,N}|t|^{-N}\norm g_2
       \quad(t\ne0,\ g\in\mathcal H).
\end{equation}
In particular, this bound is uniform over any choice of normalized
negative eigenvectors of the compressed screen. Uniformity as
\(\delta\to0\) is neither needed nor asserted.

\subsection{The fixed axis background vanishes on the scaled tests}

Local finiteness of \(\mathcal A\) implies that its nonzero
ordinate magnitudes are bounded away from zero. Choose an integer
\(N\ge1\) with \(2N>\max\{m_{\mathcal A},1\}\). Then
\begin{equation}\label{scr:inverse-moment}
 S_{\mathcal A,N}
    =\sum_{\substack{i\gamma\in\mathcal A\\\gamma\ne0}}
            \nu_{\mathcal A}(i\gamma)|\gamma|^{-2N}<\infty.
\end{equation}
Indeed the finitely many nonzero ordinates of magnitude at most
one give a finite sum. On \(2^j<|\gamma|\le2^{j+1}\),
\eqref{scr:background-count} bounds the remaining contribution
by a constant times \(2^{j(m_{\mathcal A}-2N)}\);
the resulting geometric series converges. If there are no
nonzero ordinates, the sum in \eqref{scr:inverse-moment} is zero.

For a unit vector \(g\in\mathcal H\) define
\[
 f_L(u)=L^{-1/2}g(u/L).
\]
This is a unit compact smooth test with support strictly inside
\((-L,L)\), and
\[
 H_{f_L}(i\gamma)=\sqrt L\,H_g(iL\gamma).
\]
Oddness gives \(H_{f_L}(0)=0\), so any multiplicity at the
origin contributes nothing. Equations \eqref{scr:smooth-decay}
and \eqref{scr:inverse-moment} give the uniform bound
\begin{equation}\label{scr:background-vanishing}
 0\le Q_{\mathcal A}(f_L,f_L)
   \le C_{\mathcal H,N}^2S_{\mathcal A,N}L^{1-2N}
      \longrightarrow0\qquad(L\longrightarrow\infty).
\end{equation}
No boundedness of the axis form on the full \(L^2\) space is
being assumed.

We can now prove \eqref{scr:transition}. Axis positivity and
the finite-screen eigenvalue imply
\begin{equation}\label{scr:margin-upper}
 0\le\cM_{Z_\beta}(L)\le-\mu_-(n,\beta,L).
\end{equation}
This in particular proves finiteness of \(\cM_{Z_\beta}(L)\),
even if the positive axis form is unbounded on \(L^2(-L,L)\).
At \(L=x\beta^{-q_n}\),
\[
 \epsilon=\beta L=x\beta^{1/(2d+1)},\qquad
 L\epsilon^{2d}=x^{2d+1}.
\]
For any compact interval \(I\subset(0,\infty)\),
\eqref{scr:eigenvalue} and \eqref{scr:margin-upper} therefore
give
\begin{equation}\label{scr:transition-upper}
 \cM_{Z_\beta}(x\beta^{-q_n})
    \le c_nx^{2d+1}+O_{n,I}(\beta^{2/(2d+1)})
       \quad(x\in I).
\end{equation}

For the reverse bound, take \(g=g_{\epsilon,\mathcal H}\) to
be any normalized negative eigenvector of the compressed screen.
Its scaled test has, by \eqref{scr:compressed-eigenvalue},
\[
 Q_{\mathcal S_{n,\beta}}(f_L,f_L)
   =-c_n(\mathcal H)x^{2d+1}
           +O_{n,\mathcal H,I}(\beta^{2/(2d+1)}).
\]
On \(I\), the right side of \eqref{scr:background-vanishing}
is uniformly \(O_{\mathcal H,\mathcal A,I}
(\beta^{q_n(2N-1)})\). Hence
\begin{equation}\label{scr:transition-lower}
 \cM_{Z_\beta}(x\beta^{-q_n})
   \ge c_n(\mathcal H)x^{2d+1}
                -o_{n,\mathcal H,\mathcal A,I}(1)
       \quad(x\in I).
\end{equation}
Given an error tolerance, first choose \(\mathcal H_\delta\)
so that \((c_n-c_n(\mathcal H_\delta))
\max_{x\in I}x^{2d+1}\) is smaller than that tolerance,
using \eqref{scr:coefficient-recovery}. With this smooth space
fixed, let \(\beta\to0\) in \eqref{scr:transition-upper}
and \eqref{scr:transition-lower}. This proves the locally
uniform limit \eqref{scr:transition}. The order of the two
limits avoids any assertion of uniform derivative bounds as
the cutoffs approach the endpoints.

\subsection{Inversion to the first prescribed-margin radius}

The test classes increase with \(L\), so
\(\cM_{Z_\beta}(L)\) is nondecreasing. Put
\[
 x_\eta=(\eta/c_n)^{1/(2d+1)}.
\]
For any \(0<x_-<x_\eta<x_+\), \eqref{scr:transition}
implies, for sufficiently small \(\beta\),
\[
 \cM_{Z_\beta}(x_-\beta^{-q_n})<\eta,\qquad
 \cM_{Z_\beta}(x_+\beta^{-q_n})>\eta.
\]
The first inequality, together with monotonicity, excludes the
requested margin at every radius at most \(x_-\beta^{-q_n}\).
The second supplies an actual smooth test with margin strictly
greater than \(\eta\): a supremum strictly above \(\eta\)
must have a member above \(\eta\). In particular the defining
set for \(L_\eta(\beta)\) is nonempty for small \(\beta\),
and
\[
 x_-\le\beta^{q_n}L_\eta(\beta)\le x_+.
\]
Let \(x_-\uparrow x_\eta\) and \(x_+\downarrow x_\eta\).
This proves \eqref{scr:first-radius} for both strict and
non-strict definitions, without an attainment assumption.

For \(n=1\), the formula reads
\[
 L_\eta(\beta)\sim(3\eta/4)^{1/3}\beta^{-2/3}.
\]
For \(n=2\) and \(n=3\), the exponent is \(6/7\), with the
different constants in \eqref{scr:first-coefficients}.
In general the radius exponent is \(2n/(2n+1)\) for odd
\(n\), and \((2n+2)/(2n+3)\) for even \(n\).

The fixed-background conclusion includes every fixed axis divisor
satisfying \eqref{scr:background-count}. It does not
include backgrounds whose locations or multiplicities vary with
\(\beta\). Likewise the
fixed-\(n\) expansion \eqref{scr:eigenvalue} supplies no error
bound uniform in a growing screen size. The result here is the
exact prescribed-margin transition for a fixed configuration,
not an assertion uniform in a growing count budget.
